\section*{Chapter \ref{ch:qm:covariance-estimation-and-random-matrices} --- Covariance Estimation and Random Matrices}

\begin{solution}{exo:qm:covariance-estimation-and-random-matrices:1}
$q = 0.4$: $\lambda_\pm = (1 \pm \sqrt{0.4})^2 = 0.135$ and $2.665$ (for a correlation matrix, $\sigma^2 = 1$).
\end{solution}

\begin{solution}{exo:qm:covariance-estimation-and-random-matrices:2}
$q = 0.5$, so the true volatility is about $5\%/(1 - 0.5) = 10\%$.
\end{solution}

\begin{solution}{exo:qm:covariance-estimation-and-random-matrices:3}
With $w = \Sigma^{-1}\mathbf 1/c$, $c = \mathbf 1^\top\Sigma^{-1}\mathbf 1$: $w^\top\Sigma w = \mathbf 1^\top\Sigma^{-1}\Sigma\Sigma^{-1}\mathbf 1/c^2 = c/c^2 = 1/c$.
\end{solution}

\begin{solution}{exo:qm:covariance-estimation-and-random-matrices:4}
Spikes above $1 + \sqrt{0.25} = 1.5$. A spike of 3 appears at $3(1 + 0.25/2) = 3.375$, above the edge $(1.5)^2 = 2.25$.
\end{solution}

\begin{solution}{exo:qm:covariance-estimation-and-random-matrices:5}
$\delta mI + (1 - \delta)U\Lambda U^\top = U(\delta mI + (1 - \delta)\Lambda)U^\top$. Every eigenvalue is pulled toward $m$ by the same fraction, so the small eigenvalues, which are too small,
and the large ones, which are too large, are corrected by one number; the oracle correction is not linear in $\lambda$, and \omterm{def:qm:covariance-estimation-and-random-matrices:nonlinear}{nonlinear shrinkage} chooses it eigenvalue by eigenvalue.
\end{solution}

\begin{solution}{exo:qm:covariance-estimation-and-random-matrices:6}
$1/(1 - q) \le 1.1$ requires $q \le 0.091$, so $T \ge 200/0.091 = 2\,200$ days, about 8.7 years; over such a span the covariance itself changes.
\end{solution}

\begin{solution}{exo:qm:covariance-estimation-and-random-matrices:7}
\texttt{bias\_vs\_q()}: 1.11, 1.26, 1.43, 1.67, 2.05, 2.51, 3.38, 5.21 and 10.05 at $q = 0.1$ to 0.9, against 1.11, 1.25, 1.43, 1.67, 2.00, 2.50, 3.33, 5.00 and 10.09.
\end{solution}

\begin{solution}{exo:qm:covariance-estimation-and-random-matrices:8}
The in-sample volatility of a \omterm{def:qm:covariance-estimation-and-random-matrices:pca}{minimum-variance portfolio} is lowest for the \omterm{def:qm:estimation:estimator}{estimator} whose small eigenvalues are most underestimated, that is the noisiest: in the chapter the raw
sample covariance has the lowest in-sample figure (5.7\%) and the worst realised risk (9.6\%). Select \omterm{def:qm:estimation:estimator}{estimators} by out-of-sample risk, or by the realised-to-predicted ratio.
\end{solution}

\begin{solution}{pb:qm:covariance-estimation-and-random-matrices:1}
\textbf{1.} Predicts 5.71\%, delivers 9.58\% against the truth (9.61\% over the next year); the true optimum is 7.44\%.
\textbf{2.} $1/(1 - 0.4) = 1.67$; the simulation gives 1.68.
\textbf{3.} 0.135 and 2.665; on pure noise the sample eigenvalues run from 0.144 to 2.565.
\textbf{4.} The market (51.4 against 52.9 in the population) and four sectors (2.96 to 3.68 against 2.83 to 2.99): the sector eigenvalues are pushed up by the noise, as
$\ell(1 + q/(\ell - 1))$ predicts.
\textbf{5.} A spike of $1 + \sqrt{0.4} = 1.63$ times the noise variance.
\textbf{6.} 0.034 toward the identity, delivering 9.14\%; 0.178 toward constant correlation, also 9.14\% but with an honest prediction (7.64\%).
\textbf{7.} \omterm{def:qm:covariance-estimation-and-random-matrices:nonlinear}{Nonlinear shrinkage} 8.25\% (predicts 7.36\%); clipping 8.05\% (predicts 7.27\%).
\textbf{8.} 7.74\% (predicts 6.65\%): the truth is a six-factor model, which six principal components recover, and the diagonal remainder is exactly right.
\textbf{9.} 1.68, 1.51, 1.19, 1.12, 1.11 and 1.17.
\textbf{10.} Clipping and \omterm{def:qm:covariance-estimation-and-random-matrices:nonlinear}{nonlinear shrinkage} have the most honest predictions (ratios 1.11 and 1.12); the \omterm{def:qm:covariance-estimation-and-random-matrices:factor}{factor model} builds the best portfolio (7.74\%).
\textbf{11.} Its true volatility is ten times its prediction (10.05 simulated, 10.09 by the formula).
\textbf{12.} $q \le 0.091$: 2\,200 days.
\textbf{13.} It weights each direction by the inverse of its estimated variance, and the most underestimated variances are noise; the optimiser concentrates exactly where the estimate is
worst.
\textbf{14.} It forbids the large offsetting positions that exploit noise, which acts as a form of shrinkage (at the cost of excluding genuinely useful shorts).
\textbf{15.} Stocks share a market factor; the identity target says they are uncorrelated, so a strong shrinkage toward it would be badly wrong and the optimal intensity stays tiny (0.034).
\textbf{16.} \omterm{def:qm:covariance-estimation-and-random-matrices:nonlinear}{Nonlinear shrinkage} or a \omterm{def:qm:covariance-estimation-and-random-matrices:factor}{factor model} with a shrunk remainder, cross-checked by clipping; not the raw sample.
\textbf{17.} Track realised against predicted volatility of the production portfolios on a rolling window, and the stability of the eigenvalues above the edge.
\textbf{18.} Heavier tails widen the eigenvalue spread and favour robust or rank-based estimates; a changing covariance shortens the usable history, which raises $q$ and makes cleaning
more important.
\textbf{19.} Named result: \emph{the portfolio that promised 6\%}: at $q = 0.4$ the sample \omterm{def:qm:covariance-estimation-and-random-matrices:pca}{minimum-variance portfolio} delivers 1.68 times the volatility it predicts (5.71\% promised, 9.58\%
delivered), against $1/(1 - q) = 1.67$; \omterm{def:qm:covariance-estimation-and-random-matrices:linear}{linear shrinkage} toward the identity gives 1.51, toward constant correlation 1.19, \omterm{def:qm:covariance-estimation-and-random-matrices:nonlinear}{nonlinear shrinkage} 1.12, clipping 1.11 and a six-factor model 1.17.
\textbf{20.} Because it seeks the directions of smallest estimated risk, which are the directions where estimation error has made risk look smallest.
\end{solution}

\begin{solution}{iq:qm:covariance-estimation-and-random-matrices:1}
The optimiser selects the directions whose sample variance is smallest, and those are, disproportionately, directions where noise has pushed the estimate down; in high dimension the
in-sample variance is biased down by about $1 - q$ and the true variance up by $1/(1 - q)$, $q = N/T$.
\iqlookfor{Selection on noise and the $1/(1 - q)$ scaling.}
\end{solution}

\begin{solution}{iq:qm:covariance-estimation-and-random-matrices:2}
The limiting eigenvalue density of a sample covariance of pure noise, supported on $[\sigma^2(1 - \sqrt q)^2, \sigma^2(1 + \sqrt q)^2]$. On a correlation matrix, eigenvalues above the upper edge
carry information (market, sectors); those inside the bulk are indistinguishable from noise and can be cleaned (clipped or shrunk).
\iqlookfor{The edges, and the bulk-versus-spikes reading.}
\end{solution}

\begin{solution}{iq:qm:covariance-estimation-and-random-matrices:3}
A convex combination of the sample covariance and a structured target, with the intensity that minimises the expected Frobenius loss, estimated from the data. The intensity grows with the
sampling noise of $S$ (large $N/T$, fat tails) and shrinks with the distance between $S$ and the target.
\iqlookfor{Bias-variance trade and what drives the intensity.}
\end{solution}

\begin{solution}{iq:qm:covariance-estimation-and-random-matrices:4}
$q = 1\,000/750 > 1$: the sample covariance is singular. Use a \omterm{def:qm:covariance-estimation-and-random-matrices:factor}{factor model} (statistical or fundamental) for the common part with a diagonal or shrunk remainder, or \omterm{def:qm:covariance-estimation-and-random-matrices:nonlinear}{nonlinear shrinkage}
designed for $N > T$; check by the realised risk of the portfolios it produces.
\iqlookfor{Recognising singularity and using structure.}
\end{solution}

\begin{solution}{iq:qm:covariance-estimation-and-random-matrices:5}
Those whose eigenvalues exceed the Marchenko--Pastur upper edge for the sample's $q$ (after accounting for the variance they take from the bulk): typically the market and a handful of
sectors; the rest cannot be told from noise with that much data.
\iqlookfor{The edge as the criterion, and the spike threshold.}
\end{solution}

\begin{solution}{iq:qm:covariance-estimation-and-random-matrices:6}
An \omterm{def:qm:estimation:estimator}{estimator} that keeps the sample eigenvectors and changes only the eigenvalues, so it transforms consistently with any rotation of the data. The best in the class uses the oracle
eigenvalues $u_i^\top\Sigma u_i$, the true variance of each sample eigenportfolio; \omterm{def:qm:covariance-estimation-and-random-matrices:nonlinear}{nonlinear shrinkage} estimates them from the sample spectrum through random-matrix theory.
\iqlookfor{Equivariance, the oracle, and how it is approximated.}
\end{solution}
